% ============================================
% §7 Discussion
% ============================================
\section{Discussion}
\label{sec:discussion}

\paragraph{Interpretation and scope.}
The main lesson is that the memory interface can be made into an
object of evaluation rather than left as a prompting convention. In
a closed-rule game with text-readable state, the bounded typed
contract supports fixed-$A_0$ wins, locates the largest
within-harness difference at the $L_5$ skill layer (directional at
our sample size), and keeps fixed-difficulty performance separate
from ladder endpoints. The release is meant to
make the next comparison easier: an accumulating-context variant can
be added in the same codebase, with the same condition tags, frozen
stores, prompt records, and scoring scripts, rather than inferred
from public runs that use different harnesses.

\paragraph{Implications.}
Two observations suggest broader applicability. First, separating
memory into typed slots makes attribution tractable: gains can be
traced to a specific layer rather than to ``more context.'' Second,
the bounded contract decouples interface design from accumulating
state, making the same evaluation surface portable to non-game
agentic tasks with comparable closed-rule structure.

\paragraph{Implications for loop engineering.}
The bounded contract is a concrete, measurable design point for the memory
stage of closed-rule, turn-based agent loops: per-decision typed retrieval
keeps the online context bounded regardless of run length, typed stores make
memory updates auditable, and postrun writes expose learning as explicit
artifacts rather than opaque transcript growth. Whether this pattern is the
right default for open-ended production loops is untested here; what we provide
is a reproducible testbed and archive for measuring a memory-layer change under
fixed game, denominators, and scoring---an empirical complement to the largely
qualitative loop-engineering guidance now common in
practice~\citep{anthropic-context-2025,memory-survey-2025}.
